% ============ Solutions and hints for exercises (Part 5: §13-§15) ============
% 原书 pp.292-293; 原书对 §13 的 Ex.13.1,13.2,13.5、§14 的 Ex.14.1-14.6、§15 的 Ex.15.3 未给解答, 编号按原书跳号.

\section*{\S\ref{sec:13}}

\noindent\textbf{13.3.}\quad Let $P \in \Ass(A)$ be such that $\htop P \ge 1$, and let $p_1, \ldots, p_r$ be the prime
divisors of $(a)$. If $P \neq p_i$ for all $i$, then there exists $x \in P$ such that $(a):x = (a)$.
This is a contradiction, since $x$ is a zero-divisor, but if $xy = 0$ then $y \in \sqrt{(a)A}$.

\noindent\textbf{13.4.}\quad (i) is easy. (ii) The homogeneous elements of $P$ are nilpotent mod $Q^*$,
hence so are all elements of $P$. Now we show that if $f \notin P$, $g \notin Q^*$ then
$fg \notin Q^*$. Let $f = f_1 + \cdots + f_r$, $g = g_1 + \cdots + g_s$ with $f_i$ and $g_j$ homogeneous,
and $\deg f_1 < \deg f_2 < \cdots$, $\deg g_1 < \deg g_2 < \cdots$; we work by induction on
$r + s$. If $r = s = 1$ there is nothing to prove. Also, since we can assume that
$g_1 \notin Q^*$ we have $g_1 \notin Q$. If $f_1 \notin P$ then $f_1g_1 \notin Q^*$. Next, suppose $f_1 \in P$. If
$f_1g_1 \in Q^*$ then $fg \in Q^*$, since $(f_1 + \cdots + f_r)g \in Q^*$. If $f_1g_1 \notin Q^*$ then
$f_1^tg_1 \in Q^*$ and $f_1^{t-1}g_1 \notin Q^*$ for some $t > 1$ (since $f_1^n \in Q^*$ for $n \gg 0$).
Replacing $g$ by $f_1^tg$ reduces to the case $f_1g_1 \in Q^*$, so that $f_1^tg_1 \notin Q^*$.

\noindent\textbf{13.6.}\quad First half: let $S$ be the multiplicative set made up of homogeneous elements
of $R$ not in $P$; then $R_S/P^*R_S$ can be viewed as the localisation of $R/P^*$ with
respect to all non-zero homogeneous elements, and by the previous question
this is $\cong K[X, X^{-1}]$, which is a one-dimensional ring. Second half: proof by
induction on $\htop P = n$; take a prime ideal $Q \subset P$ with $\htop Q = n - 1$. If $Q \neq P^*$
then $Q$ is inhomogeneous, $Q^* \subset P^*$ and $\htop(P/Q^*) \ge 2$, so by the first half,
$P^* \neq Q^*$, hence $\htop P^* \ge \htop Q^* + 1 = n - 1$.

\section*{\S\ref{sec:14}}

\noindent\textbf{14.7.}\quad Let $p \in \Spec A$, $f, g \notin p$, and let $l = \dim A/p$. If $l = 1$ then $pA_f$ is a maximal
ideal. Suppose that $l > 1$. We can choose $x_2, \ldots, x_r \in m$ such that
$\htop(p, f, x_2, \ldots, x_i)/p = \htop(p, g, x_2, \ldots, x_i)/p = i$ for $2 \le i \le r$. Then any mini-
mal prime divisor $P$ of $(p, x_2, \ldots, x_r)$ satisfies $\dim A/P = l$, $f \notin P$ and $g \notin P$,
so $g \notin PA_f \in \text{m-Spec}(A_f)$.

\section*{\S\ref{sec:15}}

\noindent\textbf{15.1.}\quad Set $Z = X/Y$ so that $X = YZ$, $B = k[Y, Z] \supset A = k[YZ, Y]$, and $pB =
YB$. So $B/pB \cong k[Z]$ and $\dim B_P/pB_P = 1$. Now let $p' = (X - \alpha Y)A$ for
$0 \neq \alpha \in k$; then any height $1$ prime ideal of $B$ containing $X - \alpha Y =
Y(Z - \alpha)$ must be $YB$ or $(Z - \alpha)B$, but since $YB \cap A \neq p'$ and $(Z - \alpha)B \cap A \neq p'$
there does not exist any prime ideal of $B$ contained in $P$ and lying over $p'$.

\noindent\textbf{15.2.}\quad No. Set $f = XY - 1$. Then $fB$ is a prime ideal of $B$, and $fB \cap A = (0)$.
Since $fB + XB = B$ there does not exist any prime ideal of $B$ containing
$fB$ and lying over $XA$. Note that the fibres of $A \to B$ are all one-
dimensional.
